% 两视图的对极几何；点、线位置为示意，不对应具体相机参数。宽约105mm。
\begin{tikzpicture}[x=1cm,y=1cm,font=\figurefont\small]
  \coordinate (C1) at (0,0);
  \coordinate (C2) at (7,0);
  \coordinate (X) at (3.6,4.2);
  \coordinate (Xb) at (2.52,2.94);
  % 两个像平面
  \filldraw[fill=FigureBlueFill,draw=FigureBlue,line width=0.7pt]
    (0.3,1.0) -- (2.4,1.6) -- (2.4,3.0) -- (0.3,2.4) -- cycle;
  \filldraw[fill=FigureBlueFill,draw=FigureBlue,line width=0.7pt]
    (4.4,1.4) -- (6.7,0.9) -- (6.7,2.4) -- (4.4,3.1) -- cycle;
  % 基线与对极点
  \draw[BookInk,line width=0.9pt] (C1) -- node[below] {基线} (C2);
  % 射线
  \draw[FigureRed,line width=0.9pt] (C1) -- (X);
  \draw[FigureRed,line width=0.9pt] (C2) -- (X);
  \draw[FigureYellow,dashed,line width=0.8pt] (C2) -- (Xb);
  % 像点
  \coordinate (x1) at ($(C1)!0.43!(X)$);
  \coordinate (x2) at ($(C2)!0.47!(X)$);
  \coordinate (x2b) at ($(C2)!0.52!(Xb)$);
  \fill[FigureRed] (x1) circle (1.6pt) node[left] {$\tilde{\bm x}_1$};
  \fill[FigureRed] (x2) circle (1.6pt) node[right] {$\tilde{\bm x}_2$};
  % 对极线：第二幅图中过 x2 与 x2b 的线
  \draw[FigureBlue,line width=1pt] ($(x2b)!-0.2!(x2)$) -- ($(x2b)!1.6!(x2)$);
  \node[text=FigureBlue,anchor=west,font=\figurefont\footnotesize] at (5.75,2.25) {极线 $\bm l_2=\bm F\tilde{\bm x}_1$};
  \fill[FigureRed] (X) circle (1.8pt) node[above] {$\bm X$};
  \fill[FigureYellow] (Xb) circle (1.6pt) node[above left,text=BookInk] {$\bm X'$};
  \fill[BookInk] (C1) circle (1.5pt) node[below left] {$\bm C_1$};
  \fill[BookInk] (C2) circle (1.5pt) node[below right] {$\bm C_2$};
\end{tikzpicture}
